\section{Discussion and Future Work}
\label{thesis:concl-future}

\acp{TEE} are maturing and becoming a cornerstone of modern applications. In
this dissertation, we illustrated how they can be securely and efficiently
integrated into the design of applications and protocols, but we also
highlighted some operational challenges that application developers and
administrators should take into consideration. In this section, we build upon
our contributions to reflect about possible directions for future work, focusing
more on general aspects of integrating \acp{TEE} into applications rather than
discussing specific improvements of our previous work.

\paragraph{Secure I/O on commercial TEEs}

Almost all applications require interaction with the external world to function
properly, such as end-user devices, sensors and actuators for cyber-physical
systems, and co-processors and accelerators for dedicated workloads like machine
learning. Current \acp{TEE} protect applications running in the main CPU, while
encryption in transit can protect data that goes in and out of the secure
environment. When the computation moves to another device, however, additional
isolation mechanisms are necessary.

In \cref{ch:authexec}, we described how to perform secure I/O on Sancus
microcontrollers to protect the communication between an enclave and an I/O
device such as a sensor. We achieved this by leveraging memory-mapped I/O to map
a device's memory to protected enclave memory and configuring exclusive access
of such memory to a corresponding \emph{driver} enclave
(cf.~\cref{impl:sancus}). While this solution was good enough for our small use
case, it may not be suitable for more complex I/O devices whose drivers are
proprietary and/or consist of a large number of lines of code. Furthermore, most
commercial \acp{TEE} do not support secure I/O out of the box, making our
approach unusable. A solution to this problem could be relying on proxy devices
that sit between an enclave and a I/O peripheral to encrypt and authenticate all
communication, as proposed in earlier
work~\cite{BASTION-SGX,SGX-USB,eskandarian2019fidelius}. Other approaches to
securely connect ``legacy'' devices to a \ac{TEE} would require hardware changes
and a careful design~\cite{bertschi2024devlore}. However, these solutions do not
protect the I/O devices themselves, which can still be malicious or contain
vulnerabilities. For devices that can directly access protected \ac{TEE} memory,
e.g., via \ac{DMA} operations, this poses a non-negligible risk.

A significant step towards fully secure I/O would be running a \ac{TEE} in I/O
devices and accelerators as well, which would make every computation secure
(i.e., confidentiality- and integrity-protected) and attestable in every step.
This is the approach of the Nvidia Hopper architecture~\cite{nvidiaCC}, which
allows running isolated GPU workloads that can be interconnected with
\acs{VM}-based \acp{TEE} such as Intel \ac{TDX} and AMD \sevsnp. Today, the
Hopper GPU is available on Microsoft Azure in preview mode and is already being
used for confidential computing applications such as confidential
AI~\cite{edgelessContinuumAI}. Another ongoing line of work is defining
standards to securely attach a \ac{TEE}-aware device to a secure virtual
machine, such as the \ac{TDISP}~\cite{tdispSpec}. The \ac{TDISP} standard is
currently being implemented by Intel and
AMD~\cite{intelTdxConnect,AmdSevTioWhitepaper} and there are academic solutions
for ARM \ac{CCA} as well~\cite{sridhara2024acai}.

While the solutions mentioned above seem promising, they also make management
and attestation of \ac{TEE} workloads more complex. In principle, one could
apply a similar approach as we proposed in \cref{ch:authexec}, where each
software module can run either in the main CPU (e.g., in enclaves or
confidential virtual machines), or in external devices such as a GPU. Each
module and the underlying hardware should then be attested and communication
between modules properly protected. Future work needs to investigate mechanisms
to streamline this process while keeping the attack surface as small as
possible.

\paragraph{Increase transparency in TEE software}

The security guarantees of \acp{TEE} depend on several software components, such
as the \ac{TEE} firmware and microcode, libraries and \acp{SDK}, attestation
software such as the \ac{QE} in Intel \ac{SGX}, and of course the code that runs
inside the \ac{TEE}. All of these components are part of the \ac{TCB}, and a
vulnerability in any part of the code might compromise the isolation guarantees
provided by the \ac{TEE}. Ideally, the \ac{TCB} should be as small as possible
and every software component should be verifiable to ensure that no such
vulnerabilities exist. In reality, however, the code running in a \ac{TEE} is
often large enough to make formal verification techniques unsuitable or only
usable on small parts of the codebase. The problem is exacerbated on
\acs{VM}-based \acp{TEE}, since their workloads consist of fully-fledged virtual
machines composed of several software layers, even though there are some
solutions to dramatically reduce the size of \ac{CVM} workloads, such as
Gramine-TDX~\cite{kuvaiskii2024gramine} and Mushroom~\cite{mushroomTee}. It is
essential that developers and administrators consider \ac{TCB} reduction as a
primary objective when designing a confidential application.

Furthermore, as revealed in \cref{ch:cloud}, attesting \acs{VM}-based \acp{TEE}
is a complex process and, when using confidential computing solutions in the
cloud, cloud customers have limited capabilities as some software components are
closed-source and/or directly managed by the cloud providers, making
verification harder or even impossible to achieve independently. Without proper
attestation, it is hard to tell whether the customer's data is really protected
or not, and relying on the cloud provider for attestation defeats the purpose of
using \acp{TEE} in the cloud in the first place. To make \ac{TEE} workloads
really \emph{trustworthy}, therefore, it is necessary that a verifier is capable
at least to independently verify the authenticity of all the code running inside
the \ac{TEE}. As discussed in \cref{section:cloud}, none of the major cloud
providers offers a complete solution to achieve a transparent verifiability of
the \acs{CVM} boot process, although some improvements have been made in
the past months (cf.~preamble of \cref{ch:cloud}). Future work must prioritize
this issue, focusing on open-source code and reproducible builds to increase
transparency and trustworthiness. Secure I/O solutions, as discussed above,
should also be involved in this process.

\paragraph{Overcome limitations of TEEs}

In \cref{ch:time}, we focused on the interface between an application running
inside a \ac{TEE} and the untrusted world, pointing out that some use cases
require careful design choices, especially when dealing with (untrusted) system
resources such as a clock. This problem becomes especially dangerous when
porting existing applications to \acp{TEE} with a ``lift-and-shift'' approach,
i.e., without any code changes. Since traditional applications have been
designed with a different threat model in mind, where the host is assumed to be
trusted (cf.~\cref{thesis:intro-concepts}), they might rely on resources that
might be controlled by privileged attackers. For example, the Sigy attack
exploits malicious signals to compromise unaware \ac{SGX}
applications~\cite{sridhara2024sigy}. Therefore, such issues should be
acknowledged by developers when designing a \ac{TEE} application, avoiding
lift-and-shift when possible.

Furthermore, there are some inherent limitations of \acp{TEE} that should be
taken into account. Other solutions for protecting data in use, such as
\ac{FHE}~\cite{acar2018fhe} and \ac{MPC}~\cite{lindell2020mpc}, rely on advanced
mathematics and statistics principles to transform the data before being used,
such that the computations are performed on the transformed data rather than the
original plaintext. With \acp{TEE}, instead, hardware-based isolation mechanisms
and encryption prevent unauthorized access, but the data is still processed in
the clear within the main processor. This puts significant responsibilities and
trust on \ac{TEE} manufacturers, since vulnerabilities in a \ac{TEE} might
compromise the security of applications and leak sensitive data. Unfortunately,
as highlighted in \cref{thesis:intro-threat}, this has already happened several
times in the past and is likely to happen again in the future. For example,
researchers at Google analyzed the implementation of AMD \sevsnp{} and Intel
\ac{TDX}, finding several issues that resulted in more than 30
\acp{CVE}~\cite{securingUnseen}. Besides, since the data is processed in the
clear, attackers might leverage techniques such as side-channel and physical
attacks to bypass the isolation provided by a \ac{TEE} and exfiltrate data.
Although this kind of attacks is an orthogonal problem that is partially out of
scope of the \ac{TEE} threat model (even though certain \acp{TEE} offer defense
mechanisms against some side channels), it must be taken into account when
deploying a confidential application in the real world. To mitigate such
attacks, as mentioned in \cref{thesis:intro-threat}, developers can rely on
software-based techniques such as constant-time programming and oblivious
computation: To this end, a good real-world example is the Signal contact
discovery service~\cite{signalContactDiscoveryOram}, which uses an \ac{ORAM}
algorithm to hide access patterns and protect against side-channel
attacks~\cite{stefanov2018pathoram}.

Research has shown that attacks on \acp{TEE} are a realistic threat and defenses
can only reduce the risk but cannot completely remove it. As such, future work
should investigate solutions to minimize the consequences of a successful
attack. For example, our \ac{V2X} self-revocation scheme proposed in
\cref{ch:v2x} assumes that the \ac{TC} securely stores all cryptographic
material and, in case of revocation, makes it unusable for future communication.
However, if a key leaks from the \ac{TC} due, e.g., to a side-channel attack,
there is currently no mechanism to revoke such key. Although we explicitly
clarified that these attacks are out of scope, they would partially invalidate
the properties that we formally verified
(cf.~\cref{section:design-verification}). Another interesting case study is
given by Van Schaik et al., who performed an attack on a blockchain network that
leverages Intel \ac{SGX} to provide confidentiality to smart contracts via a
secret seed shared by all enclaves~\cite{schaik2024sgxsok}. The authors were
able to recover such seed by adding a rogue \ac{SGX} enclave to the network,
compromising the confidentiality of the entire blockchain. These examples
demonstrate how important it is to consider \ac{TEE} failures in the design
phase of a real-world application, e.g., to develop ``failsafe'' mechanisms that
take effect when such failures occur. In our \ac{V2X} case, a simple solution
could be adding to our design a fixed lifetime to all certificates, similar to
the passive revocation scheme~\cite{etsi2022102941}. While most of these
mechanisms may be application-specific, some concepts could be generalized to
many use cases, such as credential management systems on the whole.

\paragraph{TEEs in industry standards}

In \cref{ch:v2x}, and in particular in \cref{section:discussion-compatibility},
we discussed how our self-revocation scheme can be integrated into existing
\ac{V2X} protocols, and we highlighted how industry standards like
ETSI~\cite{etsi2022102941} cannot support our scheme out of the box but need
some adaptations, i.e., to require and make use of a \ac{TC} in vehicles. This
showcases that, in order to increase adoption of \acp{TEE} and make them a
mainstream solution, industry standards need to integrate this technology in
their design. As a matter of fact, some standards do already exist, such as the
IETF \ac{RATS}~\cite{ietfRats} and \ac{TEEP}~\cite{ietfTeep} architectures, as
well as some reports like the ETSI NFV SEC 009~\cite{etsiNFVSecurity}. However,
they discuss \acp{TEE} in a generic way without applying them to real use cases
yet.

In order to facilitate the integration into standards, research needs to
demonstrate how \acp{TEE} can fit into modern use cases and which benefits they
provide compared to traditional solutions. Proposing a novel design is not
enough, as it needs to be accompanied by an extensive evaluation to demonstrate
its feasibility, as well as a fair comparison with the state of the art.
Additionally, formal techniques can be used to give strong arguments about the
security posture of the new design. In \ac{V2X} credential management systems,
for example, we noticed a general lack of these fundamentals, both in standards
and research, although with some exceptions~\cite{whitefield2017formal}.
